% =========================================================================
% Stochastic-FX variance proxy — working paper draft
% Repo: abrigo-analytics (HEAD 27fe058, plan v0.6 exit gate)
% Paper covers PRIMITIVES.md §15 open item 2 closure:
%   - PRIMITIVES.md §6 σ_T proxy and ε(σ_T) inversion as a unified mechanism
%   - Four model instantiations (deterministic + GBM + OU + Merton)
%   - cadCAD wrapping (bridge to dynamic-system simulation)
%   - AI-cost-factor model as the immediate downstream consumer
%   - Per-family canonical-pin verification (Phase A + Phase B + Phase C)
% Skill: /latex-econ-model (notation conventions + structure)
% =========================================================================
\documentclass[11pt,letterpaper]{article}

% --- Essential math packages (per latex-econ-model skill) ----------------
\usepackage{amsmath}
\usepackage{amssymb}
\usepackage{amsthm}
\usepackage{mathtools}
\usepackage{bm}
\usepackage{dsfont}

% --- Document hygiene -----------------------------------------------------
\usepackage[utf8]{inputenc}
\usepackage[T1]{fontenc}
% NOTE: microtype omitted (requires scalable lmodern not in this system's
% texlive-latex-recommended install). Body text uses Computer Modern at the
% default sizes; visual quality is acceptable for working-paper draft.
\usepackage[margin=1in]{geometry}
\usepackage{parskip}
\usepackage{xcolor}                           % needed for blue!50!black mix
\usepackage{hyperref}
\hypersetup{
    colorlinks=true,
    linkcolor=blue!50!black,
    urlcolor=blue!50!black,
    citecolor=blue!50!black,
}
\usepackage{booktabs}
\usepackage{array}
\usepackage{graphicx}
\graphicspath{{../../notes/stochastic_fx_figures/}}

% --- Theorem environments -------------------------------------------------
\newtheorem{definition}{Definition}[section]
\newtheorem{proposition}{Proposition}[section]
\newtheorem{lemma}{Lemma}[section]
\newtheorem{remark}{Remark}[section]
\newtheorem{assumption}{Assumption}[section]

% --- Standard econ commands -----------------------------------------------
\newcommand{\E}{\mathbb{E}}
\newcommand{\Var}{\text{Var}}
\newcommand{\Cov}{\text{Cov}}
\newcommand{\Prob}{\mathbb{P}}
\newcommand{\R}{\mathbb{R}}
\newcommand{\N}{\mathbb{N}}
\newcommand{\ind}{\mathds{1}}
\newcommand{\pd}[2]{\frac{\partial #1}{\partial #2}}
\DeclareMathOperator*{\argmax}{arg\,max}
\DeclareMathOperator*{\argmin}{arg\,min}

% --- Project-specific commands --------------------------------------------
% σ_T realized-variance proxy (PRIMITIVES.md §6 eq. 7)
\newcommand{\sigmaT}{\sigma_{T}}
% ε inversion (PRIMITIVES.md §6 eq. 8)
\newcommand{\epsT}{\varepsilon(\sigmaT)}
% FX path numerator / denominator
\newcommand{\Xbar}{\bar{X}}
\newcommand{\Ybar}{\bar{Y}}
\newcommand{\xbar}{\bar{x}}
% (Y, M, X) triple notation
\newcommand{\Ymx}{(Y,\,M,\,X)}
% Centering projection M = I - (1/N) 1 1^T
\newcommand{\centerProj}{\mathcal{M}}
% Pin tags
\newcommand{\PinZ}[1]{\textsc{pin Z1.#1}}

\title{%
    The $\sigmaT$ Variance Proxy and Its $\varepsilon$ Inversion: \\
    A Unified Mechanism Across Deterministic and Stochastic FX Dynamics \\
    \large with a cadCAD-wrapping Bridge to the AI-Cost-Factor Downstream Consumer
}
\author{%
    Juan Manuel Serrano \\
    \texttt{juan.serranotmf@gmail.com}
}
\date{2026-05-16 \quad (working paper draft)}

\begin{document}
\maketitle

\begin{abstract}
We close \textsf{notes/PRIMITIVES.md} \S15 open item 2 (\emph{stochastic-FX
variant — Path A v3}) by verifying, across three SDE families (geometric
Brownian motion, Ornstein--Uhlenbeck, and Merton jump-diffusion), that the
realized-variance proxy
\(\sigmaT = T^{-1} \sum_{j} (X_j - \bar X_{\text{path}})^2\)
admits a single family-agnostic inversion
\(\varepsilon(\sigmaT) = \sqrt{8\,\sigmaT/\xbar^{2}}\)
that holds pointwise (\PinZ{3a}, algebraic identity) and distributionally
(\PinZ{3b} mean-only moment match against literature-derived discrete-statistic
analytic moments; \PinZ{4} Kolmogorov--Smirnov goodness-of-fit against
per-family-dispatched reference distributions). Phase B (moment) and
Phase C (KS) verdicts pass at canonical pin under each SDE family at a fixed
\(N = 1000\) Monte-Carlo budget (\PinZ{5} anti-fishing floor); audit-block
hash determinism (\PinZ{2}) is bit-exact across both intra-process and
inter-process re-runs.

The unified mechanism is wrapped, in a forthcoming integration, into a
cadCAD dynamic-agent simulation where the verified $\sigmaT$ machinery
becomes the structural M-side substrate for downstream
\(\Ymx\)-triple iterations. We position the AI-cost-factor model
(\textsf{docs/specs/2026-05-16-ai-cost-factor-model-design.md} v0.2.0,
under construction) as the immediate downstream consumer: its candidate-X
identification step is independent of the verification this paper provides,
but its eventual M-design step inherits the \(\sigmaT\)-replicating
Carr--Madan strip (\textsf{cohort\_5\_strip}, deterministic instantiation)
and the per-family stochastic verification floor.

The paper makes no causal claim about the wage-to-productive-capital
ratchet that motivates the broader Abrigo operating framework; the
contribution is strictly the verification floor on the $\sigmaT$ mechanism
that any downstream \(\Ymx\) instrument family must clear before its
M-design step is admissible.
\end{abstract}

\tableofcontents

% =========================================================================
\section{Unified framework: $\sigmaT$ as a structural primitive}
\label{sec:framework}
% =========================================================================

\subsection{Operating context and the $\Ymx$-triple unit of work}
\label{sec:framework:ymx}

The Abrigo operating framework, summarized in \texttt{CLAUDE.md}, takes as
its highest-level goal the minimization of income inequality framed in
post-Keynesian terms (distribution as institutionally determined rather
than equilibrium-given). The product family pursued by the project is
\emph{permissionless on-chain perpetual convex instruments} settled on
Panoptic, with denomination chosen per target population (Mento-native
\textsf{COPm}/\textsf{BRLm}/\textsf{KESm}/\textsf{EURm}/\textsf{USDm},
USDC, ETH, or sectoral baskets). The transmission channel from wage to
productive capital runs through a \emph{premium-funded ratchet}: a wage
earner pays a small recurring premium out of wage income, and the
instrument's accumulated convex payoff and roll yield convert over time
into productive-capital exposure that absent the instrument the wage
earner would never cross the wage--capital boundary to obtain.

The operating unit of work is the \(\Ymx\)-triple:
\begin{itemize}
    \item $Y$ — outcome variable measuring a target population's exposure
        to the wage-to-capital transition;
    \item $M$ — the Panoptic pool configuration hosting the hedge
        (underlying pair, strike/range geometry, payoff shape);
    \item $X$ — the major risk currently blocking that transition for
        the target population.
\end{itemize}

This paper does not introduce a new \(\Ymx\)-triple. It verifies a
\emph{shared substrate} that every \(\Ymx\)-triple containing a convex
on-chain hedge inherits.

\subsection{The $\sigmaT$ proxy and the $\varepsilon$ inversion}
\label{sec:framework:sigmaT}

\begin{definition}[Realized-variance proxy, $\sigmaT$]
\label{def:sigmaT}
Let $\{X_j\}_{j=0}^{n}$ denote a discrete FX-rate path observed at
$n+1$ equally spaced timepoints $t_j = j \cdot \Delta t$ on the
trading horizon $[0, T]$, with $\Delta t = T/n$ and
$\bar{X}_{\text{path}} = (n+1)^{-1}\sum_{j=0}^{n} X_j$ the path's
sample mean. The discrete realized-variance proxy is
\begin{equation}
    \sigmaT \;\coloneqq\; \frac{1}{T}\sum_{j=0}^{n}
    \bigl(X_j - \bar{X}_{\text{path}}\bigr)^{2}.
    \label{eq:sigmaT}
\end{equation}
\end{definition}

Equation~\eqref{eq:sigmaT} is the discrete-time form of
\textsf{notes/PRIMITIVES.md} \S6 eq.~(7). The normalization is by $T$
(not by $n+1$): the statistic carries units of $[X]^{2}$, and it scales
proportionally to grid density. This is intentional and load-bearing:
under any continuous-time FX-path generator, the sum-of-squared
deviations sample more variation as $n \to \infty$, and dividing by $T$
preserves comparability across grids at fixed $T$.

\begin{definition}[$\varepsilon$ inversion]
\label{def:eps}
Given a realized-variance proxy $\sigmaT$ and an FX mean $\xbar$
(typically the canonical pin $\xbar = \bar{X}/\bar{Y} = 4000$ for the
COP/USD-denominated cohort), the inversion is
\begin{equation}
    \epsT \;=\; \sqrt{\frac{8\,\sigmaT}{\xbar^{2}}}.
    \label{eq:eps}
\end{equation}
\end{definition}

Equation~\eqref{eq:eps} is \textsf{PRIMITIVES.md} \S6 eq.~(8). It is
purely algebraic: no SDE-family parameter enters; the relation is
family-agnostic. The square-root form ensures (\PinZ{3a} trivial-degenerate
limit) that
\(
    \lim_{\sigmaT \to 0^{+}} \epsT \;=\; 0,
\)
which captures the geometric fact that a flat (deterministic) path has
zero implied volatility per unit FX mean.

\subsection{Verification structure: three phases, six pins}
\label{sec:framework:phases}

The verification this paper reports has three phases, each anchored to a
distinct numbered pin in
\textsf{docs/specs/2026-05-11-stochastic-fx-variant-design.md} v0.5:

\begin{description}
    \item[Phase A — algebraic identity (\PinZ{3a}).] For each path's
        $\sigmaT^{(n)}$, the pointwise identity
        $\bigl(\varepsilon(\sigmaT^{(n)})\bigr)^{2} = 8\sigmaT^{(n)} / \xbar^{2}$
        holds to float64 precision (\textsc{numerical\_identity\_tol} $= 10^{-6}$).
        Family-agnostic; reduces to a re-evaluation of~\eqref{eq:eps}.

    \item[Phase B — mean-only moment match (\PinZ{3b}, spec v0.5 amendment).]
        Empirical mean of $\sigmaT^{(n)}$ across the $N$-path ensemble
        matches the hand-derived analytic $\E[\sigmaT]$ for the
        DISCRETE statistic~\eqref{eq:sigmaT} at the canonical grid,
        within $\textsc{moment\_rel\_tol} = 0.05$. The discrete-statistic
        analytic uses the centering-projection identity
        \(
            \E[\sigmaT \cdot T] = \mathrm{tr}\bigl(\centerProj\,\Sigma\bigr)
                                  + \mu^{\top} \centerProj\, \mu
        \)
        where $\centerProj = I_{n+1} - (n+1)^{-1}\mathbf{1}\mathbf{1}^{\top}$
        and $\Sigma$ is the SDE-family-specific auto-covariance kernel
        on the canonical grid. Variance is NOT in this gate (see
        Remark~\ref{rem:mean-only}).

    \item[Phase C — Kolmogorov--Smirnov goodness-of-fit (\PinZ{4},
        spec v0.5 amendment).] Per-family dispatch on the reference
        distribution shape (not on the threshold):
        \begin{itemize}
            \item GBM: method-of-moments lognormal reference fit to the
                analytic $(\E, \Var)$, tested via
                \texttt{scipy.stats.ks\_1samp}.
            \item OU: method-of-moments gamma reference fit to the
                analytic $(\E, \Var)$, tested via \texttt{ks\_1samp}.
            \item Merton: empirical-CDF reference via a high-$N$
                ($N_{\textsc{ref}} = 100\,000$) reference run from the
                same SDE, tested via \texttt{ks\_2samp}. This avoids the
                shape mismatch between Merton's Poisson-mixture-of-lognormals
                geometry and the 2-parameter lognormal alternative.
        \end{itemize}
        Single threshold across all three families:
        $\textsc{ks\_pvalue\_floor} = 0.01$.
\end{description}

\begin{remark}[Why mean-only Phase B; what variance discloses]
\label{rem:mean-only}
At the \PinZ{5} anti-fishing path-count floor $N = 1000$, the standard
error of the sample-variance estimator on $\sigmaT$ is intrinsically
$\sqrt{(\kappa_{\text{eff}} - 1)/(N-1)}$ where $\kappa_{\text{eff}}$
ranges 4--35 across the three families (driven by jump kurtosis for
Merton). This produces an 8\%--30\% Monte-Carlo noise floor on the
variance estimator — strictly incompatible with
$\textsc{moment\_rel\_tol} = 0.05$. The spec v0.4 attempt to gate
variance against the same 5\% tolerance was mathematically
unsatisfiable; v0.5 \S11.6 dropped variance from the Phase B gate.
Variance relative error is still computed and emitted on the
\textsf{InversionVerdict} as an audit-trail observation (it is not
zero); the full-distribution match is preserved at Phase C, where the
KS test against a moment-matched reference constrains the joint shape
(mean, variance, skewness, kurtosis) via the empirical CDF.
\end{remark}

\subsection{Pin coverage and anti-fishing posture}
\label{sec:framework:pins}

The six pins anchoring this paper are:

\begin{center}
\begin{tabular}{@{}lp{0.7\textwidth}@{}}
\toprule
\textbf{Pin} & \textbf{Statement} \\
\midrule
\PinZ{1} & SDE parameters pre-pinned ex-ante per family (\S\ref{sec:gbm}, \S\ref{sec:ou}, \S\ref{sec:merton} below). \\
\PinZ{2} & Path-ensemble determinism: identical
            $(\text{params}, \text{rng\_seed}, n_{\text{paths}})$
            triples yield bit-exact paths and bit-exact audit\_block. \\
\PinZ{3a} & Algebraic identity \(\bigl(\epsT\bigr)^{2} = 8\sigmaT/\xbar^{2}\)
            to float64 precision. \\
\PinZ{3b} & Mean-only moment match against discrete-statistic analytic
            $\E[\sigmaT]$ within 5\%. \\
\PinZ{4} & Per-family Phase C KS p-value $\geq 0.01$. \\
\PinZ{5} & Anti-fishing routing: failures route to \textsc{corrections-$\alpha$}
            and scoped two-reviewer re-review; NEVER silent parameter
            re-tuning; NEVER ``increase $N$ until passing''; \(N=1000\)
            is an invariant, not a parameter to grow. \\
\PinZ{6} & Strip preservation: the
            \textsf{cohort\_5\_strip}/\textsf{IronCondor\_strip.json}
            audit\_block remains byte-equal across the verification
            plan's execution window. \\
\bottomrule
\end{tabular}
\end{center}

\PinZ{5} carries forward the broader project's anti-fishing invariants
($N_{\min} = 75$, $\textsc{power}_{\min} = 0.80$,
$\textsc{mdes\_sd} = 0.40$ SD-units of $Y$, pre-pinned sign expectation
and lag structure). For this paper specifically, three concrete
anti-fishing prohibitions hold across all four model instantiations:
(i)~no silent threshold relaxation, (ii)~no seed re-roll on failure,
(iii)~no $N$-floor lowering. The four spec amendments
(\S11.6, \S11.7 of the parent spec; \S16.3, \S16.4, \S16.5, \S16.7 of
the plan) are each subjected to scoped two-reviewer re-review per
\textsf{master-spec} \S6.4, and each preserves these three prohibitions.

\subsection{Downstream consumer: the AI-cost-factor model}
\label{sec:framework:downstream}

The verification floor established by this paper is the immediate
prerequisite for an \(\Ymx\)-triple targeting the COP-denominated AI-tooling
cost burden borne by non-subscribed LATAM developers (the AI-cost-factor
model under construction, spec v0.2.0; see
\S\ref{sec:cadcad:cost-factor}). That model's X-channel identification
step (R5 descriptive risk quantification + R4-Stage-3 vol-clustering
regression) operates on the empirical cost path, NOT on $\sigmaT$
directly. But should X-identification yield a candidate FX-vol risk
that admits Panoptic-eligible M-design, the M-side instrument is a
$\sigmaT$-replicating strip and inherits the per-family stochastic
verification floor of this paper.

This positioning is asymmetric on purpose: the $\sigmaT$ machinery is
\emph{substrate} for downstream \(\Ymx\)-triple iterations, not the
final product. The contribution of this paper is the substrate
verification; downstream consumers carry their own identification
burdens.

% =========================================================================
\section{Deterministic instantiation: Carr--Madan replication}
\label{sec:deterministic}
% =========================================================================

% -------------------------------------------------------------------------
% §2.1 — Deterministic FX-path generator (closed-form perturbation)
% -------------------------------------------------------------------------
\subsection{Deterministic FX-path generator and its $\sigmaT$ limit}
\label{sec:det:generator}

The deterministic instantiation of the framework specializes the FX
trajectory $\{X_j\}_{j=0}^{n}$ of Definition~\ref{def:sigmaT} to a closed-form
perturbation around the cohort's stationary mean
$\Xbar/\Ybar$. Following \textsf{notes/PRIMITIVES.md} \S5 eq.~(6), the
generator is
\begin{equation}
    (X/Y)_t(\varepsilon, \omega)
    \;=\;
    \Bigl(1 + \varepsilon \bigl(\cos^{2}(\omega t) - \tfrac{1}{2}\bigr)\Bigr)
    \,\frac{\Xbar}{\Ybar},
    \qquad 0 < \varepsilon < 1,
    \label{eq:det:path}
\end{equation}
with perturbation kernel $f_t \coloneqq \cos^{2}(\omega t) - \tfrac{1}{2}$ so
that $(X/Y)_t = (1 + \varepsilon f_t)\,\Xbar/\Ybar$. The amplitude
$\varepsilon$ controls the path's oscillation magnitude around the
stationary mean and the frequency $\omega$ controls its rate. The
deterministic generator is the noiseless limit of the stochastic SDE
families instantiated in \S\ref{sec:gbm}--\ref{sec:merton}: it admits a
closed-form $\sigmaT$ and so anchors the algebraic identity
$\bigl(\epsT\bigr)^{2} = 8\sigmaT/\xbar^{2}$ of equation~\eqref{eq:eps}
at $\bar X = \Xbar$, $\bar Y = \Ybar$, $\xbar = \Xbar/\Ybar$.

% -------------------------------------------------------------------------
% §2.2 — Closed-form sigma_0 anchor (R1 reduction)
% -------------------------------------------------------------------------
\subsection{The $\sigma_{0}$ closed-form anchor}
\label{sec:det:sigma0}

\begin{proposition}[Deterministic-limit reduction]
\label{prop:sigma0}
Under generator~\eqref{eq:det:path} the realized-variance proxy
$\sigmaT$ of Definition~\ref{def:sigmaT} reduces, in the long-horizon
ergodic average over the kernel $f_t$, to the closed form
\begin{equation}
    \sigma_{0}
    \;=\;
    \frac{(\Xbar/\Ybar)^{2}\,\varepsilon^{2}}{8},
    \label{eq:det:sigma0}
\end{equation}
yielding the inverse map $\varepsilon = \sqrt{8\sigma_{0}/(\Xbar/\Ybar)^{2}}$
that is the deterministic specialization of the family-agnostic
inversion~\eqref{eq:eps}.
\end{proposition}

Equation~\eqref{eq:det:sigma0} is the R1 anchor of
\textsf{notes/STAGE\_2\_RESULTS.md} \S2.1; it is obtained by substituting
$f_t = \cos^{2}(\omega t) - \tfrac{1}{2}$ into the squared-deviation sum
of equation~\eqref{eq:sigmaT} and exploiting
$\langle \cos^{2}(\omega t) - \tfrac{1}{2}\rangle_{t} = 0$,
$\langle \bigl(\cos^{2}(\omega t) - \tfrac{1}{2}\bigr)^{2}\rangle_{t} = 1/8$
on a horizon that resolves the period $2\pi/\omega$. At the canonical
cohort pin $\Xbar/\Ybar = 4000$, $\varepsilon = 0.10$, identity
\eqref{eq:det:sigma0} pins $\sigma_{0} = 2{,}000$; equation~\eqref{eq:eps}
recovers $\epsT = 0.10$ exactly.

% -------------------------------------------------------------------------
% §2.3 — Carr--Madan static replication of sigma_T
% -------------------------------------------------------------------------
\subsection{Carr--Madan static replication}
\label{sec:det:carr-madan}

The Carr--Madan static-replication theorem~\cite{carr-madan-2001}
establishes that any twice-differentiable European payoff
$g(S_{T})$ admits an exact representation as a portfolio of vanilla
out-of-the-money puts and calls, plus a bond and a forward leg:
\begin{equation}
    g(S_{T})
    \;=\;
    g(S_{0}) + g'(S_{0})(S_{T} - S_{0})
    + \int_{0}^{S_{0}} g''(K)\,(K - S_{T})^{+}\,dK
    + \int_{S_{0}}^{\infty} g''(K)\,(S_{T} - K)^{+}\,dK.
    \label{eq:carr-madan}
\end{equation}
The canonical Carr--Madan variance contract is the log-contract,
$g(S_{T}) = -2\log(S_{T}/S_{0}) + 2(S_{T} - S_{0})/S_{0}$, for which
$g''(K) = 2/K^{2} > 0$ — the canonical convex weight profile that, when
integrated against the option-strike continuum~\eqref{eq:carr-madan},
returns $\sigmaT$ at expiry up to the usual martingale correction.

\begin{remark}[From integral to discrete strip]
\label{rem:strip-discretization}
The continuous representation~\eqref{eq:carr-madan} is not Panoptic-eligible
as written; Panoptic's pool mechanics admit only finite-leg long-vol
positions on Uniswap v3/v4 ranges. The deployment-eligible
discretization replaces the two strike integrals in~\eqref{eq:carr-madan}
by a finite sum of IronCondor cells, each cell contributing a piecewise-linear
convex bump on a strike window. The 12-leg three-condor strip introduced
in \S\ref{sec:det:strip} is the smallest Panoptic-feasible discretization
that preserves long-vol convexity throughout the convex span.
\end{remark}

% -------------------------------------------------------------------------
% §2.4 — Twelve-leg IronCondor strip geometry
% -------------------------------------------------------------------------
\subsection{The 12-leg IronCondor strip geometry}
\label{sec:det:strip}

The replicating instrument is a three-condor strip stacked at progressively
wider strikes around the cohort's stationary spot $S_{0} = \Xbar/\Ybar$.
Each condor contributes four legs in the long-vol convention
(\emph{reverse} IronCondor: short the outermost strikes, long the inner
strikes):
\begin{equation}
    P_{\text{cond}}(S_{T}; K_{1},K_{2},K_{3},K_{4})
    \;=\;
    -(K_{1} - S_{T})^{+} + (K_{2} - S_{T})^{+}
    + (S_{T} - K_{3})^{+} - (S_{T} - K_{4})^{+},
    \label{eq:condor-payoff}
\end{equation}
with strikes $K_{1} < K_{2} < K_{3} < K_{4}$. The strip payoff is the
weighted sum
\begin{equation}
    \Pi_{\text{strip}}(S_{T})
    \;=\;
    \sum_{j=1}^{3} w_{j}\,
    P_{\text{cond}}\bigl(S_{T};\,K_{1}^{(j)},K_{2}^{(j)},K_{3}^{(j)},K_{4}^{(j)}\bigr),
    \label{eq:strip-payoff}
\end{equation}
with three condor cells labelled $\textsf{left\_tail}$, $\textsf{atm}$,
$\textsf{right\_tail}$ in the committed strip artefact and weights $w_{j}$
solved as the least-squares projection of $\Pi_{\text{strip}}$ onto the
log-contract proxy. The tiled-body geometry, where the inner body of each
condor abuts its neighbour's center, drives the strip's floor at spot
exactly to zero — $\Pi_{\text{strip}}(S_{0}) = 0$ within $10^{-9}\cdot S_{0}$
— so that the long-vol payoff signature
\begin{equation}
    \Pi_{\text{strip}}(S_{0}\,e^{-\delta_{\text{inner}}}) > 0
    \quad\text{and}\quad
    \Pi_{\text{strip}}(S_{0}\,e^{+\delta_{\text{inner}}}) > 0
    \label{eq:long-vol-signature}
\end{equation}
is satisfied strictly on both sides of spot (Pin S5 in
\textsf{cohort\_5\_strip}). The strict-positivity requirement on both
wings is what the long-vol predicate enforces; this is the
\textsc{corrections-$\alpha$} amendment from the Stage-3 A1 disposition,
which deprecated the v0.2 strategy survey's permissive zero-floor
predicate.

\begin{lemma}[Log-contract envelope tracking]
\label{lem:envelope}
Let $Q(S_{T}) \coloneqq -2\log(S_{T}/S_{0}) + 2(S_{T} - S_{0})/S_{0}$ be
the Carr--Madan log-contract proxy and let
$\Pi_{\text{strip}}^{c}(S_{T}) \coloneqq \Pi_{\text{strip}}(S_{T}) -
\Pi_{\text{strip}}(S_{0})$ be the spot-centered strip payoff. On a symmetric
log-strike grid spanning $[S_{0} e^{-\delta}, S_{0} e^{+\delta}]$ with
$\delta = 0.15$, the best-fit scale $\beta^{\star} =
\sum_{i} \Pi_{\text{strip}}^{c}(S_{i})\,Q(S_{i})\big/\sum_{i} Q(S_{i})^{2}$
yields max relative residual
\(
    \max_{i} \bigl|\Pi_{\text{strip}}^{c}(S_{i})
                     - \beta^{\star}\,Q(S_{i})\bigr| \big/
    \max_{i} \bigl|\Pi_{\text{strip}}^{c}(S_{i})\bigr|
    \;\leq\;
    0.05,
\)
the canonical 5\% envelope tolerance of \textsf{PRIMITIVES.md} \S12.
\end{lemma}

Lemma~\ref{lem:envelope} formalizes Pin S6 of the cohort-5 strip:
the spot-centered strip tracks the log-contract proxy up to a single
best-fit scale, with residual within the same 5\% band that
\PinZ{3b} reserves for moment matching of the stochastic instantiations.
The constant payoff floor absorbed by centering is held by the LP as
static collateral; the convex \emph{excess} above that floor is what the
$\sigmaT$-replicating leg pays.

% -------------------------------------------------------------------------
% §2.5 — Selection rationale and Pin Z1.6 anchor
% -------------------------------------------------------------------------
\subsection{Selection rationale and Pin Z1.6 anchor}
\label{sec:det:selection}

The IronCondor strip is not asserted to be the unique discretization of
\eqref{eq:carr-madan}. It is the best-available Panoptic-eligible
long-vol strategy that emerged from the Stage-3 A1 20-primitive survey,
in which three long-vol challengers (\textsf{long\_straddle},
\textsf{zeehbs}, \textsf{long\_strangle}) underperformed the IronCondor
baseline by 52, 64, and 69 percentage points respectively on the
envelope-residual metric of Lemma~\ref{lem:envelope}. Selection is
locked under verdict \textsc{keep\_ironcondor}.

The strip artefact is committed at
\textsf{simulations/saas\_builder/data/IronCondor\_strip.json} with
\texttt{audit\_block = 94150326332b90e50cfe02b580e6d05280100b430de0089ea9197c8fa4aaf329}.
\PinZ{6} preserves this hash byte-equal across the verification window
of this paper: any downstream stochastic instantiation that consumes the
strip as its deterministic limit reads it through this hash anchor.

\paragraph{Implementation pointer.}
The geometry construction lives in
\textsf{simulations/saas\_builder/cohort\_5\_strip/geometry.py}; the
envelope verifier and the long-vol signature assertion in
\textsf{cohort\_5\_strip/replication.py}; the strip-with-audit-block emit
in \textsf{cohort\_5\_strip/emit.py}. The Stage-2 cohort-5 work was
completed prior to this paper's framework and is cited here, not
re-derived, as the deterministic limit of the unified $\sigmaT$ mechanism.

% =========================================================================
\section{Stochastic instantiation I: geometric Brownian motion}
\label{sec:gbm}
% =========================================================================

% -------------------------------------------------------------------------
% §3.1 — GBM SDE and Itô log-scale discretization
% -------------------------------------------------------------------------
\subsection{The GBM SDE and its log-scale discretization}
\label{sec:gbm:sde}

The first stochastic instantiation of the framework lifts the
deterministic generator~\eqref{eq:det:path} to a geometric Brownian
motion~\cite{hull-2017}: the FX-rate process $\{X_t\}_{t \in [0,T]}$
solves the linear Itô SDE
\begin{equation}
    dX_{t} \;=\; \mu\,X_{t}\,dt + \sigma\,X_{t}\,dW_{t},
    \qquad X_{0} = x_{0} > 0,
    \label{eq:gbm:sde}
\end{equation}
where $W_{t}$ is a standard Brownian motion under the physical measure
and $(\mu, \sigma) \in \R \times \R_{>0}$ are the drift and volatility
parameters. Equation~\eqref{eq:gbm:sde} admits the explicit
log-multiplicative solution
$X_{t} = x_{0}\,\exp\bigl((\mu - \tfrac{1}{2}\sigma^{2})\,t + \sigma\,W_{t}\bigr)$,
where the $-\tfrac{1}{2}\sigma^{2}$ Itô correction is load-bearing for
the unbiased-mean property $\E[X_{t}] = x_{0}\,e^{\mu t}$.

\begin{proposition}[Euler--Maruyama log-scale discretization]
\label{prop:gbm:disc}
For an equally spaced grid $t_{j} = j\,\Delta t$ with $\Delta t = T/n$,
$j = 0, 1, \dots, n$, the exact log-scale recursion for~\eqref{eq:gbm:sde}
is
\begin{equation}
    \log X_{j+1} - \log X_{j}
    \;=\;
    \bigl(\mu - \tfrac{1}{2}\sigma^{2}\bigr)\,\Delta t
    \;+\; \sigma\sqrt{\Delta t}\,Z_{j+1},
    \qquad Z_{j+1} \stackrel{\text{iid}}{\sim} \mathcal{N}(0, 1),
    \label{eq:gbm:euler}
\end{equation}
which simulates the GBM path without time-discretization error in the
log scale.
\end{proposition}

The Itô correction $-\tfrac{1}{2}\sigma^{2}$ in~\eqref{eq:gbm:euler} is
not a discretization artefact but a structural feature: omitting it
biases the simulated $\E[X_{j}]$ by a factor of $e^{j\sigma^{2}\Delta t / 2}$,
which at the canonical pin $\sigma^{2}T = 0.01$ corresponds to a $0.5\%$
bias at the terminal point — outside the Phase B tolerance band of
\PinZ{3b} even before any auto-covariance term enters.

% -------------------------------------------------------------------------
% §3.2 — Auto-covariance kernel and centering-projection identity
% -------------------------------------------------------------------------
\subsection{Auto-covariance kernel and discrete-statistic moments}
\label{sec:gbm:autocov}

The auto-covariance kernel of the log-multiplicative solution
of~\eqref{eq:gbm:sde} at $\mu = 0$ is
\begin{equation}
    \Cov(X_{j}, X_{k})
    \;=\;
    x_{0}^{2}\bigl(\exp(\sigma^{2}\,\min(t_{j}, t_{k})) - 1\bigr),
    \qquad j, k = 0, 1, \dots, n,
    \label{eq:gbm:cov}
\end{equation}
which follows directly from $\E[X_{j}X_{k}] = x_{0}^{2}\,\phi_{2}^{\min(j,k)}$
with $\phi_{2} = \E[e^{2L}] = e^{\sigma^{2}\Delta t}$, $\phi_{1} = \E[e^{L}] = 1$
under the martingale pin. The discrete-statistic
analytic of equation~\eqref{eq:sigmaT} follows from the centering-projection
identity announced in \S\ref{sec:framework:phases}:

\begin{lemma}[Centering projection for $\sigmaT \cdot T$]
\label{lem:gbm:proj}
Let $\centerProj \coloneqq I_{n+1} - (n+1)^{-1}\mathbf{1}\mathbf{1}^{\top}$
be the centering projection on $\R^{n+1}$, and let $\Sigma$ be the
$(n+1)\times(n+1)$ matrix with entries~\eqref{eq:gbm:cov}, $\mu \in \R^{n+1}$
the mean vector $\mu_{j} = \E[X_{j}] = x_{0}$ for $\mu = 0$ canonical.
Then
\begin{equation}
    \E[\sigmaT \cdot T]
    \;=\;
    \mathrm{tr}\bigl(\centerProj\,\Sigma\bigr)
    \;+\; \mu^{\top}\centerProj\,\mu,
    \label{eq:gbm:trace-id}
\end{equation}
where the second term vanishes identically because
$\mu_{j} \equiv x_{0}$ is constant in $j$ and $\centerProj\,\mathbf{1} = 0$.
\end{lemma}

The trace identity~\eqref{eq:gbm:trace-id} reduces the discrete-statistic
$\E[\sigmaT]$ to a closed sum over the GBM covariance kernel and is the
load-bearing analytic for Phase B at this family. The continuous-time
limit (cf. \textsf{notes/stochastic\_fx\_tex/sigma\_t\_moments\_gbm.tex},
transcribed below as a sanity check on the literature-anchored form) is
\begin{equation}
    \E[\sigmaT]_{\text{cont}}
    \;=\;
    \frac{\Xbar^{2}}{\Ybar^{2}}
    \left(-1 + \frac{e^{T\sigma^{2}} - 1}{T\sigma^{2}}\right),
    \qquad
    \Var[\sigmaT]_{\text{cont}}
    \;=\;
    \frac{2\,\Xbar^{4}\,(e^{T\sigma^{2}} - 1)^{2}}{T^{2}\,\Ybar^{4}\,\sigma^{4}},
    \label{eq:gbm:cont}
\end{equation}
with the trailing $-1$ inside the bracket required for the
$\sigma \to 0$ vanishing limit (a correction logged in
\textsf{notes/stochastic\_fx\_tex/sigma\_t\_moments\_gbm.tex} under
\textsc{new-block-mq-4}). The discrete value~\eqref{eq:gbm:trace-id}
differs from~\eqref{eq:gbm:cont} by the resolution of the squared-deviation
sum at finite $n$; at the canonical grid this discrete-statistic analytic
evaluates to $\E[\sigmaT] = 1.117130\times 10^{7}$, against which the
$N = 1000$ Monte-Carlo ensemble is gated in Phase B.

\begin{remark}[Why the Isserlis Var formula is audit-trail only]
\label{rem:gbm:var}
The same projection identity yields a leading-order Isserlis-Gaussian
variance
$\Var[\sigmaT \cdot T] = 2\,\mathrm{tr}((\centerProj\,\Sigma)^{2}) +
4\,(\centerProj\,\mu)^{\top}\Sigma\,(\centerProj\,\mu)$, which is exact
only under joint-Gaussianity of $X$. Because each $X_{j}$ is lognormal
under GBM, the formula is leading-order in $\sigma^{2}T$; at the
canonical pin $\sigma^{2}T = 0.01$ the residual against the true
lognormal variance is approximately $8\%$ — strictly incompatible with
the $5\%$ moment tolerance. Spec v0.5 \S11.6 amends the Phase B gate to
the mean only (Remark~\ref{rem:mean-only} above); the variance relative
error remains on the verdict record as an audit-trail observation.
\end{remark}

% -------------------------------------------------------------------------
% §3.3 — Phase A / B / C canonical-pin verdict
% -------------------------------------------------------------------------
\subsection{Three-phase verdict at canonical pin}
\label{sec:gbm:verdict}

The GBM verification at the canonical pin
$(\mu, \sigma, x_{0}, T) = (0,\, 0.10/\sqrt{12},\, 4000,\, 12)$ with
$n = 5000$ Euler steps and $N = 1000$ paths (\PinZ{5} N-floor exact)
returns the following triple, anchored to the committed verdict
artefact:

\begin{description}
    \item[Phase A (\PinZ{3a}, algebraic identity).]
        Maximum residual of
        $\bigl(\epsT^{(n)}\bigr)^{2} - 8\sigmaT^{(n)}/\xbar^{2}$
        across the $N = 1000$ paths is $7.105\times 10^{-15}$,
        well below the $\textsc{numerical\_identity\_tol} = 10^{-6}$
        floor. The identity is family-agnostic and reduces to
        re-evaluating~\eqref{eq:eps}.
    \item[Phase B (\PinZ{3b}, mean-only moment match).]
        Empirical mean of $\sigmaT^{(n)}$ matches the analytic value
        $\E[\sigmaT] = 1.117130\times 10^{7}$ of
        equation~\eqref{eq:gbm:trace-id} with relative error
        $1.121\times 10^{-2}$, inside the $5\%$ band of
        $\textsc{moment\_rel\_tol}$. Variance relative error is
        $2.127\times 10^{-1}$, consistent with Remark~\ref{rem:gbm:var};
        it is recorded as audit trail, not gated.
    \item[Phase C (\PinZ{4}, KS goodness-of-fit).]
        Per-family dispatch fits a method-of-moments lognormal
        reference to the analytic pair $(\E[\sigmaT], \Var[\sigmaT])$
        via the standard moment-matching relations
        $s = \sqrt{\log\bigl(1 + V/E^{2}\bigr)}$,
        $\mathrm{scale} = E / \sqrt{1 + V/E^{2}}$, and tests via
        \texttt{scipy.stats.ks\_1samp}. The Kolmogorov--Smirnov
        $p$-value is $1.229\times 10^{-1}$, comfortably above the
        $\textsc{ks\_pvalue\_floor} = 0.01$.
\end{description}

Composite verdict at canonical pin: \textsc{pass} on all three phases;
$\textsf{composite\_pass} = \mathrm{true}$ by the AND-invariant of the
\textsf{InversionVerdict} Value type.

\paragraph{Why a lognormal reference for Phase C.}
The Phase C dispatch on GBM uses a moment-matched two-parameter
lognormal as the reference distribution, not the empirical CDF nor a
Gaussian. The motivation is geometric: each $X_{j}$ is itself
lognormal under~\eqref{eq:gbm:sde}, and the quadratic form
$\sigmaT \cdot T = X^{\top}\centerProj\,X$ is dominated, in the
small-$\sigma^{2}T$ regime that the canonical pin inhabits
($\sigma^{2}T = 0.01$), by a mixture of squared-lognormal moments
whose first two moments coincide with those of a single fitted
lognormal up to fourth-order terms. Among the two-parameter
alternatives that admit closed-form moment matching, the lognormal
is the only one that vanishes at zero (consistent with
$\sigmaT \geq 0$) AND has the right tail asymmetry. The OU and
Merton dispatches, by contrast, route to gamma and empirical-CDF
references respectively (see \S\ref{sec:ou}, \S\ref{sec:merton}); the
$\textsc{ks\_pvalue\_floor} = 0.01$ threshold is invariant across the
three references — only the reference shape is family-dispatched,
which is consistent with \PinZ{5} (no threshold tuning).

\begin{remark}[Cross-process determinism]
\label{rem:gbm:audit}
The verdict artefact at
\textsf{simulations/stochastic\_fx/data/inversion\_verdict\_gbm.json}
carries
\(
    \texttt{audit\_block} = \texttt{3a0f75401d517f6d}\ldots
\)
(64-char lowercase hex). Identical
$(\text{params}, \text{rng\_seed}, n_{\text{paths}})$ triples replayed
from a fresh checkout reproduce this prefix bit-exact, discharging
\PinZ{2}.
\end{remark}

\paragraph{Implementation pointer.}
The path generator lives in
\textsf{simulations/stochastic\_fx/generators.py}
(class \texttt{GBMPathGenerator}, Euler--Maruyama log-scale per
Proposition~\ref{prop:gbm:disc}); the continuous-time moments
in \textsf{simulations/stochastic\_fx/moments.py}
(\texttt{gbm\_sigma\_t\_moments}); the discrete-statistic analytic
in \textsf{simulations/stochastic\_fx/variance\_proxy.py}
(\texttt{gbm\_discrete\_sigma\_t\_moments}); and the canonical pin in
\textsf{simulations/stochastic\_fx/types.py}
(\texttt{CANONICAL\_GBM}). The hand-derivation script
\textsf{scratch/2026-05-13-task-4.2-discrete-moments/derivation.py}
cross-checks the closed form of Lemma~\ref{lem:gbm:proj} against an
independent Monte-Carlo run; the LaTeX fragment is committed at
\textsf{notes/stochastic\_fx\_tex/sigma\_t\_moments\_gbm.tex} and
transcribed inline into~\eqref{eq:gbm:cont} above.

% =========================================================================
\section{Stochastic instantiation II: Ornstein--Uhlenbeck}
\label{sec:ou}
% =========================================================================

% -------------------------------------------------------------------------
% §4.1 — OU SDE and the EXACT Vasicek conditional-Gaussian transition
% -------------------------------------------------------------------------
\subsection{The OU SDE and its exact Vasicek transition}
\label{sec:ou:sde}

The second stochastic instantiation of the framework replaces the
log-multiplicative GBM dynamics of~\eqref{eq:gbm:sde} by a mean-reverting
Ornstein--Uhlenbeck process~\cite{vasicek-1977, karatzas-shreve-1991}: the
FX-rate process $\{X_{t}\}_{t \in [0, T]}$ solves the linear Gaussian
Itô SDE
\begin{equation}
    dX_{t} \;=\; \theta\,\bigl(\bar\mu - X_{t}\bigr)\,dt
                 \;+\; \sigma\,dW_{t},
    \qquad X_{0} = x_{0} > 0,
    \label{eq:ou:sde}
\end{equation}
with mean-reversion speed $\theta > 0$, long-run level
$\bar\mu > 0$, and diffusion volatility $\sigma > 0$. Unlike~\eqref{eq:gbm:sde}
the OU drift is additive (not multiplicative) and pulls $X_{t}$ back toward
$\bar\mu$ at exponential rate $\theta$; the family is the canonical
mean-reverting Gaussian process and admits closed-form stationary moments.

\begin{proposition}[Exact Vasicek conditional-Gaussian transition]
\label{prop:ou:disc}
For an equally spaced grid $t_{j} = j\,\Delta t$ with $\Delta t = T/n$, the
one-step transition density of~\eqref{eq:ou:sde} is exactly Gaussian: with
$Z_{j+1} \stackrel{\text{iid}}{\sim} \mathcal{N}(0, 1)$,
\begin{equation}
    X_{j+1}
    \;=\;
    \bar\mu \;+\; \bigl(X_{j} - \bar\mu\bigr)\,e^{-\theta\,\Delta t}
    \;+\;
    \sigma\,\sqrt{\,\frac{1 - e^{-2\theta\,\Delta t}}{2\theta}\,}\,Z_{j+1}.
    \label{eq:ou:vasicek}
\end{equation}
This is NOT an Euler--Maruyama discretization: it is the exact
conditional-Gaussian transition law of the OU SDE on the grid, with no
time-discretization error.
\end{proposition}

Proposition~\ref{prop:ou:disc} is the load-bearing distinction between the
OU and GBM discretizations of this paper. Euler--Maruyama applied
to~\eqref{eq:ou:sde} would carry an $O(\Delta t^{1/2})$ strong-convergence
residual; the closed-form transition~\eqref{eq:ou:vasicek} carries
\emph{none}. At the canonical pin (Definition~\ref{def:sigmaT};
$\theta = 1$, $\bar\mu = x_{0} = 4000$, $\sigma = 0.10 \cdot 4000/\sqrt{2}$,
$T = 12$, $n = 5000$), the stationary variance of $X_{t}$ equals
$\sigma^{2}/(2\theta) = 4 \times 10^{4}$ and the mean-reversion half-life
is $\log 2 / \theta \approx 0.69$, so the canonical horizon resolves
$\approx 17$ relaxation times — the process is squarely in its stationary
regime.

% -------------------------------------------------------------------------
% §4.2 — Auto-covariance kernel and centering-projection identity at x_0 = mu_bar
% -------------------------------------------------------------------------
\subsection{Auto-covariance kernel and discrete-statistic moments}
\label{sec:ou:autocov}

Iterating~\eqref{eq:ou:vasicek} from the deterministic initial condition
$X_{0} = x_{0}$ yields the closed-form auto-covariance kernel
\begin{equation}
    \Cov(X_{j}, X_{k})
    \;=\;
    \frac{\sigma^{2}}{2\theta}\,
    e^{-\theta\,|t_{k} - t_{j}|}\,
    \Bigl(1 - e^{-2\theta\,\min(t_{j}, t_{k})}\Bigr),
    \qquad j, k = 0, 1, \dots, n,
    \label{eq:ou:cov}
\end{equation}
which is the standard Vasicek kernel of \cite{karatzas-shreve-1991}~\S5.6.
The mean dynamics are $\E[X_{j}] = \bar\mu + (x_{0} - \bar\mu)\,e^{-\theta t_{j}}$;
at the canonical pin $x_{0} = \bar\mu$ the second term vanishes and the mean
becomes constant in $j$.

\begin{lemma}[Centering projection for $\sigmaT \cdot T$, canonical OU pin]
\label{lem:ou:proj}
Let $\centerProj \coloneqq I_{n+1} - (n+1)^{-1}\mathbf{1}\mathbf{1}^{\top}$,
let $\Sigma$ be the $(n+1)\times(n+1)$ matrix with entries~\eqref{eq:ou:cov},
and let $\mu \in \R^{n+1}$ be the mean vector. Then
\begin{equation}
    \E[\sigmaT \cdot T]
    \;=\;
    \mathrm{tr}\bigl(\centerProj\,\Sigma\bigr)
    \;+\; \mu^{\top}\centerProj\,\mu,
    \label{eq:ou:trace-id}
\end{equation}
and at the canonical pin $x_{0} = \bar\mu$ the second term vanishes
identically because $\mu_{j} \equiv \bar\mu$ is constant in $j$ and
$\centerProj\,\mathbf{1} = 0$, leaving the closed reduction
$\E[\sigmaT \cdot T] = \mathrm{tr}(\centerProj\,\Sigma)$.
\end{lemma}

The centering-projection identity~\eqref{eq:ou:trace-id} is the same
load-bearing analytic that drives Lemma~\ref{lem:gbm:proj} for GBM; the
two families differ only in the auto-covariance kernel substituted into
$\Sigma$. The continuous-time stationary closed forms
(\textsf{notes/stochastic\_fx\_tex/sigma\_t\_moments\_ou.tex}) collapse
to
\begin{equation}
    \E[\sigmaT]_{\text{cont, stat}}
    \;=\;
    \frac{\sigma^{2}}{2\theta},
    \qquad
    \Var[\sigmaT]_{\text{cont, stat}}
    \;=\;
    \frac{\sigma^{4}}{2\,T\,\theta^{2}},
    \label{eq:ou:cont-stat}
\end{equation}
which at the canonical pin pins $\E[\sigmaT] = 4 \times 10^{4}$ exactly.
This is the unique family in this paper where the continuous-time
stationary $\E[\sigmaT]$ admits a one-line closed form independent of $T$
— a direct consequence of OU stationarity. The discrete-statistic
analytic~\eqref{eq:ou:trace-id} differs from~\eqref{eq:ou:cont-stat} only
by the $(1 - e^{-2\theta\,\min(t_{j}, t_{k})})$ transient factor in
$\Sigma$, which is negligible on the canonical grid that resolves
$\approx 17$ relaxation times.

% -------------------------------------------------------------------------
% §4.3 — Phase A / B / C verdict at canonical pin
% -------------------------------------------------------------------------
\subsection{Three-phase verdict at canonical pin}
\label{sec:ou:verdict}

The OU verification at the canonical pin returns the following triple,
anchored to the committed verdict artefact:

\begin{description}
    \item[Phase A (\PinZ{3a}, algebraic identity).]
        Maximum residual of
        $\bigl(\epsT^{(n)}\bigr)^{2} - 8\sigmaT^{(n)}/\xbar^{2}$
        across the $N = 1000$ paths is $3.553\times 10^{-15}$,
        well below the $\textsc{numerical\_identity\_tol} = 10^{-6}$
        floor. As at GBM this is a re-evaluation of~\eqref{eq:eps} and is
        family-agnostic.
    \item[Phase B (\PinZ{3b}, mean-only moment match).]
        Empirical mean of $\sigmaT^{(n)}$ matches the discrete-statistic
        analytic value from~\eqref{eq:ou:trace-id} with relative error
        $1.681\times 10^{-3}$, comfortably inside the $5\%$ band of
        $\textsc{moment\_rel\_tol}$. Variance relative error is
        $2.711\times 10^{-2}$, consistent with the Isserlis Gaussian
        quadratic form being EXACT under OU (no truncation residual,
        only MC sampling noise); it is recorded as audit trail, not
        gated.
    \item[Phase C (\PinZ{4}, KS goodness-of-fit).]
        Per-family dispatch fits a method-of-moments gamma reference
        to the analytic pair $(\E[\sigmaT], \Var[\sigmaT])$ via the
        moment-matching relations
        \(
            a = E^{2}/V,\quad \mathrm{scale} = V/E,
        \)
        and tests via \texttt{scipy.stats.ks\_1samp} against the
        gamma CDF. The Kolmogorov--Smirnov $p$-value is
        $3.033\times 10^{-1}$, comfortably above the
        $\textsc{ks\_pvalue\_floor} = 0.01$.
\end{description}

Composite verdict at canonical pin: \textsc{pass} on all three phases;
$\textsf{composite\_pass} = \mathrm{true}$ by the AND-invariant of the
\textsf{InversionVerdict} Value type.

\paragraph{Why a gamma reference for Phase C.}
The Phase C dispatch on OU uses a moment-matched two-parameter gamma as
the reference distribution. The motivation is geometric: under~\eqref{eq:ou:sde}
the joint law of $\{X_{j}\}_{j=0}^{n}$ is multivariate Gaussian, and
$\sigmaT \cdot T = X^{\top}\centerProj\,X$ is a positive-semidefinite
quadratic form in jointly Gaussian variables. The exact distribution of
such a quadratic form is a weighted sum of independent
$\chi^{2}_{1}$ random variables (the eigen-decomposition of $\centerProj\,\Sigma$);
the moment-matched gamma is the natural two-parameter family approximating
this weighted-$\chi^{2}$ object with the same first two moments. Among
positive-support alternatives, gamma is the conjugate choice for
Gaussian-quadratic forms — strictly preferable to lognormal (chosen for
GBM, where the underlying $X_{j}$ themselves are lognormal). Per
\PinZ{5}, only the reference shape is family-dispatched; the
$\textsc{ks\_pvalue\_floor} = 0.01$ threshold is invariant across all
three references.

\begin{remark}[Mean-only Phase B Type-I rate at OU]
\label{rem:ou:type-i}
The spec v0.5 \S11.6 disclosure logs an empirical Phase B Type-I rate of
$\approx 0\%$ for OU at $N = 1000$ (against the seed=42 canonical pin
that returned $0.168\%$ mean rel-err). This is the lowest Type-I rate
across the three SDE families: OU's $\sigmaT$ has the smallest effective
kurtosis ($\kappa_{\text{eff}} \approx 4$ from the Gaussian-quadratic form),
so the standard-error floor $\sqrt{(\kappa_{\text{eff}} - 1)/(N - 1)}$ of
Remark~\ref{rem:mean-only} is $\approx 5.5\%$ at $N = 1000$ — well clear
of the $5\%$ tolerance for the mean-only gate that v0.5 retained. By
contrast, the Merton family (\S\ref{sec:merton}) inhabits the opposite
end of the kurtosis surface, with a Type-I rate $\approx 30\%$ that
\emph{would} have invalidated the mean-only gate had OU been the only
benchmark.
\end{remark}

\begin{remark}[Cross-process determinism]
\label{rem:ou:audit}
The verdict artefact at
\textsf{simulations/stochastic\_fx/data/inversion\_verdict\_ou.json}
carries
\(
    \texttt{audit\_block} = \texttt{0589f35b10f1ca7f}\ldots
\)
(64-char lowercase hex). Identical
$(\text{params}, \text{rng\_seed}, n_{\text{paths}})$ triples replayed
from a fresh checkout reproduce this prefix bit-exact, discharging
\PinZ{2}.
\end{remark}

\paragraph{Implementation pointer.}
The path generator lives in
\textsf{simulations/stochastic\_fx/generators.py}
(class \texttt{OUPathGenerator}, exact Vasicek transition per
Proposition~\ref{prop:ou:disc}); the continuous-time stationary moments
in \textsf{simulations/stochastic\_fx/moments.py}
(\texttt{ou\_sigma\_t\_moments}); the discrete-statistic analytic
in \textsf{simulations/stochastic\_fx/variance\_proxy.py}
(\texttt{ou\_discrete\_sigma\_t\_moments}, which assembles the
$\Sigma$-matrix from~\eqref{eq:ou:cov} and applies the centering
projection of~\eqref{eq:ou:trace-id}); and the canonical pin in
\textsf{simulations/stochastic\_fx/types.py}
(\texttt{CANONICAL\_OU}). The LaTeX fragment of~\eqref{eq:ou:cont-stat}
is committed at
\textsf{notes/stochastic\_fx\_tex/sigma\_t\_moments\_ou.tex}.

% =========================================================================
\section{Stochastic instantiation III: Merton jump-diffusion}
\label{sec:merton}
% =========================================================================

% -------------------------------------------------------------------------
% §5.1 — Merton SDE and its log-scale discretization with conditional-Gaussian
%       compound-Poisson aggregation
% -------------------------------------------------------------------------
\subsection{The Merton SDE and its conditional-Gaussian aggregation}
\label{sec:merton:sde}

The third and final stochastic instantiation of the framework extends
the GBM dynamics of~\eqref{eq:gbm:sde} by a compound-Poisson jump
component~\cite{merton-1976, andersen-piterbarg-2010}: the FX-rate
process $\{X_{t}\}_{t \in [0, T]}$ solves the jump-diffusion SDE
\begin{equation}
    dX_{t}
    \;=\;
    \mu\,X_{t}\,dt
    \;+\; \sigma\,X_{t}\,dW_{t}
    \;+\; X_{t^{-}}\,(J - 1)\,dN_{t},
    \qquad X_{0} = x_{0} > 0,
    \label{eq:merton:sde}
\end{equation}
where $N_{t}$ is a Poisson process of intensity $\lambda \geq 0$
independent of the Brownian motion $W_{t}$, and $\log J \sim
\mathcal{N}(\mu_{J}, \sigma_{J}^{2})$ is the lognormal jump multiplier
with $\sigma_{J} > 0$. Setting $\lambda = 0$ collapses~\eqref{eq:merton:sde}
to the GBM SDE~\eqref{eq:gbm:sde}; non-zero $\lambda$ injects
heavy-tailed Poisson-mixture-of-lognormals geometry into the marginal
laws of $\{X_{t}\}$.

\begin{proposition}[Conditional-Gaussian aggregated jump per step]
\label{prop:merton:disc}
For an equally spaced grid $t_{j} = j\,\Delta t$, the one-step log-increment
$\log X_{j+1} - \log X_{j}$ decomposes as a diffusion log-step plus a
\emph{compound-Poisson aggregated} log-jump:
\begin{align}
    \log X_{j+1} - \log X_{j}
    \;&=\;
    \underbrace{\bigl(\mu - \tfrac{1}{2}\sigma^{2}\bigr)\,\Delta t
                + \sigma\sqrt{\Delta t}\,Z_{j+1}}_{\text{diffusion step}}
    \notag\\
    \;&\qquad+\;
    \underbrace{N_{j+1}\,\mu_{J}
                + \sqrt{N_{j+1}}\,\sigma_{J}\,Z'_{j+1}}_{\text{aggregated jump step}},
    \label{eq:merton:euler}
\end{align}
where $N_{j+1} \sim \mathrm{Poisson}(\lambda\,\Delta t)$ is the per-cell
jump count and $Z_{j+1}, Z'_{j+1} \stackrel{\text{iid}}{\sim} \mathcal{N}(0, 1)$
are independent diffusion / conditional-jump standard normals. Conditional
on the count $N_{j+1} = N$, the sum of $N$ independent jump log-multipliers
$\sum_{i=1}^{N}\log J_{i}$ is exactly distributed as
$\mathcal{N}(N\mu_{J}, N\sigma_{J}^{2})$, so~\eqref{eq:merton:euler}
represents the entire compound-Poisson contribution by a single
conditional-Gaussian draw.
\end{proposition}

The aggregation identity in Proposition~\ref{prop:merton:disc} is the
load-bearing computational economy of the Merton instantiation: the naive
implementation would sample $N_{j+1}$ individual lognormal multipliers per
cell, costing $O(\text{total jumps})$ work; the conditional-Gaussian
aggregation costs $O(n_{\text{paths}} \cdot n_{\text{steps}})$ work
\emph{independent of $\lambda$}, and is exact in distribution (not an
approximation).

% -------------------------------------------------------------------------
% §5.2 — Auto-covariance kernel with jump-MGF building blocks; non-constant drift
% -------------------------------------------------------------------------
\subsection{Auto-covariance kernel and non-constant drift}
\label{sec:merton:autocov}

The Merton path-mean and auto-covariance kernel are most cleanly stated in
terms of two per-step characteristic-function building blocks
(Andersen--Piterbarg \cite{andersen-piterbarg-2010}~\S2.7):
\begin{equation}
    \log\varphi_{1}
    \;=\; \mu\,\Delta t
        + \lambda\,\Delta t\,\bigl(e^{\mu_{J} + \sigma_{J}^{2}/2} - 1\bigr),
    \qquad
    \log\varphi_{2}
    \;=\; 2\mu\,\Delta t + \sigma^{2}\,\Delta t
        + \lambda\,\Delta t\,\bigl(e^{2\mu_{J} + 2\sigma_{J}^{2}} - 1\bigr),
    \label{eq:merton:phi}
\end{equation}
so that the path mean is $\E[X_{j}] = x_{0}\,\varphi_{1}^{j}$ and the
second moment is $\E[X_{j} X_{k}] = x_{0}^{2}\,\varphi_{2}^{\min(j,k)}\,
\varphi_{1}^{|k-j|}$. The auto-covariance kernel follows by subtraction:
\begin{equation}
    \Cov(X_{j}, X_{k})
    \;=\;
    x_{0}^{2}\,\varphi_{1}^{|k-j|}\,
    \Bigl(\varphi_{2}^{\min(j, k)} - \varphi_{1}^{2\,\min(j, k)}\Bigr),
    \qquad j, k = 0, 1, \dots, n.
    \label{eq:merton:cov}
\end{equation}

\begin{remark}[Non-constant drift; full $\mu^{\top}\centerProj\,\mu$ expansion required]
\label{rem:merton:drift}
Unlike the GBM canonical pin ($\mu = 0 \Rightarrow \varphi_{1} = 1$
$\Rightarrow \mu_{j} \equiv x_{0}$) and the OU canonical pin
($x_{0} = \bar\mu \Rightarrow \mu_{j} \equiv \bar\mu$), the Merton
canonical pin has $\mu = 0, \mu_{J} = 0$, but the jump-MGF factor
$e^{\sigma_{J}^{2}/2} - 1 \neq 0$ inflates the per-step mean to
$\log\varphi_{1} = \lambda\,\Delta t\,(e^{\sigma_{J}^{2}/2} - 1) > 0$.
Hence $\mu_{j} = x_{0}\,\varphi_{1}^{j}$ is geometrically increasing in
$j$; the second term $\mu^{\top}\centerProj\,\mu$
in~\eqref{eq:gbm:trace-id} does NOT vanish and must be retained in full.
The discrete-statistic analytic
$\E[\sigmaT \cdot T] = \mathrm{tr}(\centerProj\,\Sigma) +
\mu^{\top}\centerProj\,\mu$ accordingly admits no GBM/OU-style mean-term
collapse for Merton.
\end{remark}

The remark above is the structural reason the discrete-statistic analytic
for Merton requires hand-derivation of \emph{both} trace terms; the
implementation in
\textsf{simulations/stochastic\_fx/variance\_proxy.py} assembles
$\mu$ and $\Sigma$ from~\eqref{eq:merton:phi}--\eqref{eq:merton:cov} and
applies the full centering-projection identity~\eqref{eq:gbm:trace-id}
without simplification. The continuous-time closed forms
(\textsf{notes/stochastic\_fx\_tex/sigma\_t\_moments\_merton.tex}) sum
the diffusion-plus-jump contributions in expectation
\begin{align}
    \E[\sigmaT]_{\text{cont}}
    \;&=\;
    \frac{\Xbar^{2}}{\Ybar^{2}}\,T\,\sigma^{2}
    \;+\;
    \frac{\Xbar^{2}}{\Ybar^{2}}\,T\,\lambda\,
    \bigl(e^{2\mu_{J} + 2\sigma_{J}^{2}}
          - 2\,e^{\mu_{J} + \sigma_{J}^{2}/2} + 1\bigr),
    \label{eq:merton:cont-mean}\\
    \Var[\sigmaT]_{\text{cont}}
    \;&=\;
    \frac{2\,\Xbar^{4}}{\Ybar^{4}}\,T\,\sigma^{4}
    \;+\;
    \frac{\Xbar^{4}}{\Ybar^{4}}\,T\,\lambda\,M_{4},
    \label{eq:merton:cont-var}
\end{align}
where the compound-Poisson fourth-jump moment $M_{4} = \E[(e^{J} - 1)^{4}]$
expands via the Gaussian MGF identity $\E[e^{kJ}] = \exp(k\mu_{J} +
k^{2}\sigma_{J}^{2}/2)$. At the canonical pin
$(\mu, \sigma, \lambda, \mu_{J}, \sigma_{J}, x_{0}, T) =
(0,\, 0.05/\sqrt{12},\, 1,\, 0,\, 0.10,\, 4000,\, 12)$,
equation~\eqref{eq:merton:cont-mean} pins $\E[\sigmaT]_{\text{cont}} =
1.994\times 10^{6}$, which the discrete-statistic
analytic~\eqref{eq:gbm:trace-id} reproduces within the
$N = 1000$ Monte-Carlo noise floor.

% -------------------------------------------------------------------------
% §5.3 — Phase A / B / C verdict; the Phase C BLOCK that triggered v0.4 → v0.5
% -------------------------------------------------------------------------
\subsection{Three-phase verdict and the empirical-CDF dispatch}
\label{sec:merton:verdict}

The Merton verification at the canonical pin returns the following triple:

\begin{description}
    \item[Phase A (\PinZ{3a}, algebraic identity).]
        Maximum residual of
        $\bigl(\epsT^{(n)}\bigr)^{2} - 8\sigmaT^{(n)}/\xbar^{2}$
        across the $N = 1000$ paths is $2.274\times 10^{-13}$, well below
        the $\textsc{numerical\_identity\_tol} = 10^{-6}$ floor. The two
        orders-of-magnitude inflation over GBM ($7.105\times 10^{-15}$)
        and OU ($3.553\times 10^{-15}$) reflects the wider numeric range
        of Merton $\sigmaT$ samples (jump kurtosis pushes terminal-cell
        sample variance up by $\approx 10^{6}$ over the diffusion-only
        families) — the identity remains family-agnostic.
    \item[Phase B (\PinZ{3b}, mean-only moment match).]
        Empirical mean of $\sigmaT^{(n)}$ matches the discrete-statistic
        analytic with relative error $7.015\times 10^{-3}$, comparable
        to GBM ($1.121\times 10^{-2}$) and OU ($1.681\times 10^{-3}$)
        despite Merton's substantially higher $\kappa_{\text{eff}}$.
        Variance relative error is $2.727$, consistent with
        Remark~\ref{rem:merton:var-rel-err}; it is recorded as audit
        trail, not gated.
    \item[Phase C (\PinZ{4}, KS goodness-of-fit).]
        Per-family dispatch fits an \emph{empirical-CDF} reference via
        a high-$N$ reference run (see Remark~\ref{rem:merton:ecdf}) and
        tests via \texttt{scipy.stats.ks\_2samp}. The Kolmogorov--Smirnov
        $p$-value is $1.160\times 10^{-1}$, comfortably above the
        $\textsc{ks\_pvalue\_floor} = 0.01$.
\end{description}

Composite verdict at canonical pin: \textsc{pass} on all three phases;
$\textsf{composite\_pass} = \mathrm{true}$ by the AND-invariant of the
\textsf{InversionVerdict} Value type.

\begin{remark}[Why variance relative error inflates to $2.727$ at Merton]
\label{rem:merton:var-rel-err}
The Isserlis Gaussian-quadratic-form variance formula
$\Var[\sigmaT \cdot T] = 2\,\mathrm{tr}((\centerProj\,\Sigma)^{2}) +
4\,(\centerProj\,\mu)^{\top}\Sigma\,(\centerProj\,\mu)$ is exact only
under joint Gaussianity of $\{X_{j}\}$. Under~\eqref{eq:merton:sde}
the marginal laws of $X_{j}$ are Poisson-mixtures of lognormals; the
Gaussian formula recovers only the diffusion-and-second-moment portion
of $\Var[\sigmaT]$ and structurally under-estimates the fourth-cumulant
contribution by $\approx 71\%$. The audit-trail relative error of
$2.727$ is consistent with this structural under-estimate. Spec v0.5
\S11.6 retains the mean-only Phase B gate (Remark~\ref{rem:mean-only}
above); the joint shape match across higher moments is preserved at
Phase C via the empirical-CDF dispatch below.
\end{remark}

\begin{remark}[The Phase C BLOCK and the v0.4 $\to$ v0.5 amendment]
\label{rem:merton:ecdf}
Under spec v0.4 the Phase C dispatch was a moment-matched two-parameter
lognormal reference (the GBM choice, generalized to all three families).
Application of that reference to the Merton canonical pin returned a
KS $p$-value of $3.41\times 10^{-21}$, $19$ orders of magnitude below the
$0.01$ floor. The empirical shape of Merton $\sigmaT$ at the canonical
pin has skewness $10.4$ and excess kurtosis $174$; a two-parameter
lognormal carries skewness $\leq 6.18$ at the moment-matched scale and
cannot represent Poisson-mixture-of-lognormals geometry by shape, only
by location-plus-scale.

Per spec v0.5 \S11.7, this Phase C BLOCK routed to \textsc{corrections-$\alpha$}
and the user disposition (2026-05-13) was a per-family dispatch on the
reference \emph{shape}, with the threshold left invariant per \PinZ{5}:
Merton's reference becomes an empirical-CDF drawn from a high-$N$
auxiliary run of the same \textsf{JumpDiffusionPathGenerator} under
pinned constants $N_{\textsc{ref}} = 100{,}000$ and
$N_{\textsc{ref-seed}} = 20260513$, tested via
\texttt{scipy.stats.ks\_2samp} against the $N = 1000$ tested ensemble.
The five concrete anti-fishing prohibitions of the disposition are
preserved: (i)~single threshold across all three families
($\textsc{ks\_pvalue\_floor} = 0.01$ invariant); (ii)~$N_{\textsc{ref}}$
PINNED, not implementer-tunable; (iii)~$N_{\textsc{ref-seed}}$ PINNED,
not implementer-tunable; (iv)~the reference is a DIFFERENT
sample from the same SDE, NOT \texttt{.fit()}-against-tested-sample
(NIT-MQ-1 tautology preserved); (v)~the interpretation of \PinZ{4}
(``per-family KS goodness-of-fit at $p \geq 0.01$'') is unchanged in
substance. The Phase C $p$-value of $1.160\times 10^{-1}$ reported above
is from the canonical run at hardware-constrained $N_{\textsc{ref}} =
5000$ per FLAG-RC-V0.5-1; the production constant $N_{\textsc{ref}} =
100{,}000$ produces an equivalent KS-$p$ in the $0.12$--$0.14$
ballpark per the Wave-2 robustness probe.
\end{remark}

\begin{remark}[Phase B Type-I rate $\approx 30\%$; seed-42 luck disclosure]
\label{rem:merton:type-i}
Spec v0.5 \S11.6 logs an empirical Phase B Type-I rate of $\approx 30\%$
for Merton at $N = 1000$, against the seed=42 canonical pin that
returned $0.70\%$ mean rel-err. The standard-error floor
$\sqrt{(\kappa_{\text{eff}} - 1)/(N - 1)}$ of Remark~\ref{rem:mean-only}
evaluates to $\approx 19\%$ for Merton at $N = 1000$ — well above the
$5\%$ tolerance for the mean-only Phase B gate. The seed=42 PASS is
therefore a Type-I-favorable observation, not a structural property of
the family; the spec disposition is to disclose the rate honestly rather
than tighten the gate (which would require $N \gtrsim 30{,}000$ to bring
the standard-error floor below $5\%$, in violation of the \PinZ{5}
$N$-floor invariant). Combined with the Phase C $\approx 1\%$ Type-I rate
(Smirnov 1948), the Merton family carries the highest joint Type-I
budget of the three SDEs in this paper; \PinZ{5} requires routing any
failure on a re-run of the canonical pin to \textsc{corrections-$\alpha$}
(re-derive analytic or refine discretization), NEVER to seed re-roll.
\end{remark}

\begin{remark}[Cross-process determinism]
\label{rem:merton:audit}
The verdict artefact at
\textsf{simulations/stochastic\_fx/data/inversion\_verdict\_merton.json}
carries
\(
    \texttt{audit\_block} = \texttt{12890ce95a36ffea}\ldots
\)
(64-char lowercase hex). Identical
$(\text{params}, \text{rng\_seed}, n_{\text{paths}})$ triples replayed
from a fresh checkout reproduce this prefix bit-exact, discharging
\PinZ{2}.
\end{remark}

\paragraph{Implementation pointer.}
The path generator lives in
\textsf{simulations/stochastic\_fx/generators.py}
(class \texttt{JumpDiffusionPathGenerator}, Euler--Maruyama log-scale
diffusion with conditional-Gaussian aggregated jump per
Proposition~\ref{prop:merton:disc}); the RNG-call order
(diffusion-$Z$ $\to$ Poisson counts $\to$ conditional-jump-$Z'$) is
load-bearing for \PinZ{2} and reproduced verbatim in the generator's
class docstring. The continuous-time moments live in
\textsf{simulations/stochastic\_fx/moments.py}
(\texttt{merton\_sigma\_t\_moments}); the discrete-statistic analytic
in \textsf{simulations/stochastic\_fx/variance\_proxy.py}
(\texttt{merton\_discrete\_sigma\_t\_moments}, with full
non-constant-$\mu_{j}$ expansion per Remark~\ref{rem:merton:drift});
the empirical-CDF reference assembly in the same file
(\texttt{\_merton\_reference\_sigma\_t}); and the canonical pin in
\textsf{simulations/stochastic\_fx/types.py}
(\texttt{CANONICAL\_MERTON}). The hand-derivation script
\textsf{scratch/2026-05-13-task-4.2-discrete-moments/derivation.py}
cross-checks the closed form against an independent Monte-Carlo run;
the LaTeX fragment of~\eqref{eq:merton:cont-mean}--\eqref{eq:merton:cont-var}
is committed at
\textsf{notes/stochastic\_fx\_tex/sigma\_t\_moments\_merton.tex}.

% =========================================================================
\section{cadCAD wrapping: bridge to dynamic-system simulation}
\label{sec:cadcad}
% =========================================================================

\subsection{Motivation and scope}
\label{sec:cadcad:motivation}

\paragraph{Section charge for content fill.}
Subsection explains why a cadCAD wrapping is necessary even after the
$\sigmaT$ machinery verifies under three SDE families: the verification
floor establishes that the inversion mechanism holds, but downstream
\(\Ymx\)-triple iterations need to simulate the interaction of multiple
state variables (FX path, cost path, replicating-portfolio position,
hedge accumulation) over time. The cadCAD framework
(\textsf{cadCAD-org/cadCAD}) provides a dynamic-agent simulation engine
with state variables, policy functions, and partial state update blocks
(PSUBs) suitable for this composition.

\subsection{State variables and PSUB design}
\label{sec:cadcad:psub}

\paragraph{Section charge for content fill.}
Subsection introduces the cadCAD state-variable signature: at each time
step $t$, the simulation tracks the FX path $X_t$ (sampled from one of
the three SDE families plus the deterministic baseline), the cost
trajectory $C_t$ (for the AI-cost-factor model, downstream consumer),
the replicating-portfolio's $\sigmaT$ accumulation, the wage-funded
premium contribution rate, and the convex-payoff-to-productive-capital
ratchet variable. PSUBs:
(i)~FX-path-advance (SDE-family-dispatched);
(ii)~$\sigmaT$ accumulator update;
(iii)~premium-payment-and-position-update;
(iv)~ratchet-conversion at threshold crossings.

\subsection{The AI-cost-factor model as downstream consumer}
\label{sec:cadcad:cost-factor}

The AI-cost-factor model specified at
\textsf{docs/specs/2026-05-16-ai-cost-factor-model-design.md} v0.2.0
(currently under construction; v0.1.0 rejected by all three reviewers
with five \textsc{block}s on identification and unitization grounds)
targets the COP-denominated AI-tooling cost burden borne by
non-subscribed LATAM developers. The cost-factor's identification stage
(\S2.1 R5 descriptive risk quantification + \S2.2 R4-Stage-3 vol-clustering
regression on log-COP-cost-returns) is independent of the
verification this paper provides; PASS on the cost-factor's
identification opens an M-design step that consumes the verified
$\sigmaT$-replicating mechanism.

The coupling in the cadCAD simulation is asymmetric:
\begin{enumerate}
    \item The FX-path PSUB (one of the four SDE families) produces $X_t$.
    \item The cost-trajectory PSUB consumes $X_t$ and emits $C_t$ as a
        function of (token usage demand, API rate-card, $X_t$).
    \item The $\sigmaT$ accumulator (PRIMITIVES.md \S6 eq.~(7) discretely)
        consumes the FX-path window and emits the realized variance.
    \item The replicating-portfolio PSUB (Carr--Madan 12-leg IronCondor
        instantiation) consumes $\sigmaT$ and the inversion
        $\epsT$ to construct the hedge position.
    \item The wage-funded premium PSUB drains a fixed share of the
        cohort's wage income each step, recapitalizing the position.
    \item The ratchet variable accumulates the convex payoff and crosses
        the productive-capital-equivalence threshold per the cohort's
        denomination.
\end{enumerate}

This paper does not implement the cadCAD wrapping; the implementation
work is the immediate next deliverable, specified to follow doc-first
sequencing from this paper's framework. The cost-factor model's R5/R4-S3
results, when they land, will be reported in a companion paper.

\subsection{What the cadCAD wrapping does NOT do}
\label{sec:cadcad:nondoes}

\paragraph{Section charge for content fill.}
Subsection lists the explicit exclusions: cadCAD wrapping does NOT
re-validate the $\sigmaT$ machinery (that's this paper's job); does NOT
estimate the cost-factor model's identification parameters (that's the
cost-factor spec's job); does NOT cross the wage-to-capital boundary
empirically (the ratchet variable's threshold crossings are simulated
under the cohort's assumed denomination, not empirically observed).

% =========================================================================
\section{Results: canonical-pin verification per family}
\label{sec:results}
% =========================================================================

\paragraph{Section charge for content fill.}
Section reports the canonical-pin verification verdict per family,
sourced from \textsf{notes/STOCHASTIC\_FX\_RESULTS.md} (the
authoritative Stage-3 open-item 2 closure document). Cross-family
summary table; per-family $(\sigma, \theta, \lambda, \sigma_J)$ parameter
pins; per-family $(\text{Phase A residual}, \text{Phase B mean rel-err}, \text{Phase B var rel-err (audit-trail)}, \text{Phase C KS p-value}, \text{audit\_block})$
tuples. Inline cross-check via \textsf{scratch/2026-05-11-stochastic-fx-execution/driver.py}
re-runnable from a fresh checkout (\PinZ{2} cross-process determinism).
Per-family interpretation paragraphs: GBM (lognormal $\sigmaT$ shape
absorbs both first and second moments); OU (Gaussian quadratic form
yields exact $\E[\sigmaT] = \sigma^{2}/(2\theta)$ at $x_0 = \bar\mu$;
gamma reference fits cleanly); Merton (mixture geometry requires the
empirical-CDF Phase C dispatch per spec v0.5 \S11.7).

% =========================================================================
\section{Discussion: what the verification floor unlocks}
\label{sec:discussion}
% =========================================================================

\paragraph{Section charge for content fill.}
Section synthesizes the four model-instantiation results into a single
statement of what the verification floor unlocks for downstream
\(\Ymx\)-triple iterations. Three honest limitations: (i)~mean-only
Phase B at $N = 1000$ has a 0\%--30\% Type-I rate per family
(spec v0.5 \S11.6 disclosure); (ii)~the Merton Phase C reference at
hardware-constrained $N_{\textsc{ref}} = 5000$ matches the production
$N_{\textsc{ref}} = 100\,000$ in the same KS p-value ballpark per the
Wave-2 robustness probe (KS p $\approx 0.12$--$0.14$), but production
hardware $\geq 8$ GB free is the binding constraint for the canonical
reference; (iii)~the deterministic eq.~(6) generator produces synchronized
paths whose ergodic content is artificial, motivating the stochastic
extension this paper closes. Forward statement: AI-cost-factor model
(downstream consumer) inherits this floor; cadCAD wrapping is the next
deliverable; an instrument-deployment study under live Panoptic
liquidity is gated on M-design step success in a future \(\Ymx\)-iteration
not yet specified.

% =========================================================================
\appendix
% =========================================================================

\section{Per-family analytic-moment derivations}
\label{app:moments}

\paragraph{Appendix charge for content fill.}
Per-family derivation of $\E[\sigmaT]$ for the DISCRETE statistic
\eqref{eq:sigmaT} via the projection identity
$\E[\sigmaT \cdot T] = \mathrm{tr}(\centerProj\,\Sigma) + \mu^{\top}\centerProj\,\mu$.
Sub-appendices A.1 GBM, A.2 OU, A.3 Merton. Each cites the textbook
auto-covariance kernel + transcribes the centering-projection
computation. The scratch derivation script
\textsf{scratch/2026-05-13-task-4.2-discrete-moments/derivation.py} is
the load-bearing artefact.

\section{Anti-fishing audit trail across spec amendments}
\label{app:antifishing}

\paragraph{Appendix charge for content fill.}
Appendix walks the four spec amendments (\S11.6, \S11.7 of the parent
spec; \S16.3, \S16.4, \S16.5, \S16.7 of the plan) and confirms each
preserves the three concrete anti-fishing prohibitions (no silent
threshold relaxation, no seed re-roll, no $N$-floor lowering). Each
amendment's Wave-1 RC + MQ Delphi-independent re-review verdict files
are cited; Wave-2 RC + MQ at plan-exit are cited.

\section{cadCAD-implementation spec (forward-looking)}
\label{app:cadcad-spec}

\paragraph{Appendix charge for content fill.}
Appendix sketches the cadCAD implementation spec that this paper
nominates for the immediate next deliverable: directory structure
(\textsf{simulations/cadcad/}), state-variable declarations, PSUB
implementations per the six PSUB roles itemized in
\S\ref{sec:cadcad:psub}, integration tests crossing the four FX
families and the AI-cost-factor cost trajectory, anti-fishing
invariants carried forward.

% =========================================================================
% References
% =========================================================================
\bibliographystyle{plain}
% TODO: produce notes/papers-bib/stochastic-fx-paper.bib with Hull,
%       Karatzas-Shreve, Andersen-Piterbarg, Carr-Madan, Merton (1976),
%       Vasicek (1977), Itô, cadCAD whitepaper, Panoptic whitepaper,
%       post-Keynesian distribution literature.
\begin{thebibliography}{10}
\bibitem{hull-2017}
    Hull, John C. (2017). \emph{Options, Futures, and Other Derivatives},
    10th ed. Pearson. \S15.

\bibitem{karatzas-shreve-1991}
    Karatzas, Ioannis \& Steven E. Shreve (1991).
    \emph{Brownian Motion and Stochastic Calculus},
    2nd ed. Springer. \S5.6 (Ornstein--Uhlenbeck process).

\bibitem{andersen-piterbarg-2010}
    Andersen, Leif B. G. \& Vladimir V. Piterbarg (2010).
    \emph{Interest Rate Modeling, Volume I: Foundations and Vanilla Models}.
    Atlantic Financial Press. \S2.7 (jump-diffusion moments).

\bibitem{carr-madan-2001}
    Carr, Peter \& Dilip Madan (2001).
    Optimal Positioning in Derivative Securities.
    \emph{Quantitative Finance} 1(1), 19--37.

\bibitem{merton-1976}
    Merton, Robert C. (1976).
    Option Pricing When Underlying Stock Returns Are Discontinuous.
    \emph{Journal of Financial Economics} 3, 125--144.

\bibitem{vasicek-1977}
    Vasicek, Oldrich (1977).
    An Equilibrium Characterization of the Term Structure.
    \emph{Journal of Financial Economics} 5(2), 177--188.

\bibitem{cadcad}
    cadCAD: complex Adaptive Dynamics Computer-Aided Design.
    \url{https://github.com/cadCAD-org/cadCAD}.

\bibitem{panoptic-2023}
    Panoptic: Perpetual Options on Uniswap.
    Technical whitepaper. 2023.

\bibitem{abrigo-stochastic-fx-spec}
    Abrigo Analytics (2026). Stochastic-FX variant design spec, v0.5.
    \texttt{docs/specs/2026-05-11-stochastic-fx-variant-design.md},
    in \texttt{abrigo-analytics} repository.

\bibitem{abrigo-stochastic-fx-plan}
    Abrigo Analytics (2026). Stochastic-FX variant implementation plan, v0.6.
    \texttt{docs/plans/2026-05-11-stochastic-fx-variant.md},
    in \texttt{abrigo-analytics} repository.

\bibitem{abrigo-ai-cost-factor}
    Abrigo Analytics (2026). AI-Cost Factor Model design spec, v0.2.0
    (under construction).
    \texttt{docs/specs/2026-05-16-ai-cost-factor-model-design.md},
    in \texttt{abrigo-analytics} repository.

\end{thebibliography}

\end{document}
